\section{Seleção simultânea da árvore e dos átomos}
\label{sec:variable-dp}

Para uma célula $J$ e um inteiro $d\ge0$, seja
\[
 \Representations_{d,H}(J;u,v)
\]
o conjunto das representações contínuas enraizadas em $J$, com profundidade
relativa no máximo $d$, folhas rotuladas por $\Library_H$ e valores externos
$u,v$. Diferentemente da seção anterior, a forma da árvore não é fixada.

A recorrência usa o princípio de otimalidade de \citet{bellman1957}. A
aproximação de curvas por segmentos via programação dinâmica já aparece em
\citet{bellman1961}, com antecedentes para segmentos lineares em
\citet{stone1961} e desenvolvimentos posteriores para aproximação poligonal
ótima em \citet{perezvidal1994}; aqui o estado contém os dois valores de
fronteira e a decisão de manter ou subdividir uma célula diádica.

\begin{definition}[recorrência de Bellman com árvore variável]
\label{def:variable-recurrence}
Para $d=0$, ponha
\begin{equation}
 C_0(J;u,v)=\Phi_H(J;u,v).
 \label{eq:variable-base}
\end{equation}
Para $d\ge1$, se $J_L,J_R$ são os filhos de $J$, defina
\begin{align}
 C_d(J;u,v)
 =\min\Bigl\{&
 \Phi_H(J;u,v),\notag\\
 &\lambda_{\rm bits}+\lambda_{\rm split}
 +\min_{w\in\Values_{D,H}(K)}
 \bigl[C_{d-1}(J_L;u,w)+C_{d-1}(J_R;w,v)\bigr]
 \Bigr\}.
 \label{eq:variable-recurrence}
\end{align}
\end{definition}

O primeiro ramo mantém $J$ como folha; o segundo subdivide, escolhe o valor
racional intermediário e reutiliza os ótimos dos dois filhos. A profundidade
global $D$ é usada somente para fixar um alfabeto comum suficientemente
grande; na chamada em $J$, $d$ é o número de níveis ainda disponíveis.

\begin{theorem}[exatidão, existência e determinismo]
\label{thm:variable-correctness}
Seja $J$ uma célula descendente de $K$ na profundidade relativa $r$. Para todo
$d\le D-r$ e todo $(u,v)\in\Values_{D,H}(K)^2$,
\begin{equation}
 C_d(J;u,v)
 =\min_{R\in\Representations_{d,H}(J;u,v)}
 \cost_{\mathcal Z}(R).
 \label{eq:variable-correctness}
\end{equation}
Consequentemente,
\begin{equation}
 C_{D,H}^{\star}(K)
 :=\min_{u,v\in\Values_{D,H}(K)}C_D(K;u,v)
 =\min_{R\in\Representations_{D,H}^{\Library}(K)}
 \cost_{\mathcal Z}(R).
 \label{eq:variable-global}
\end{equation}
O mínimo é finito quando
$H\ge\max\{2,k_0+1,|m_0|+1\}$. Com o desempate
\eqref{eq:tie-key}, o algoritmo devolve uma única representação, independentemente
da ordem em que estados, átomos ou valores intermediários sejam gerados.
\end{theorem}

\begin{proof}
A prova é por indução em $d$. Para $d=0$, apenas a árvore-folha é permitida,
e \eqref{eq:variable-base} coincide com seu custo ou vale $+\infty$ quando o
par é inadmissível.

Suponha o resultado para $d-1$. Toda representação de profundidade no máximo
$d$ está em exatamente um de dois casos: a raiz é folha; ou a raiz é interna
e possui duas subárvores de profundidade no máximo $d-1$. No segundo caso, a
continuidade produz um único valor compartilhado $w$. Reciprocamente, a
colagem de quaisquer duas subárvores com o mesmo $w$ é admissível. A
decomposição do custo é a mesma do \cref{thm:fixed-correctness}; portanto a
partição exaustiva dos candidatos nos dois casos é exatamente
\eqref{eq:variable-recurrence}. A hipótese de indução completa
\eqref{eq:variable-correctness}.

Minimizar os estados externos na raiz prova \eqref{eq:variable-global}. Sob a
hipótese sobre $H$, o átomo de parâmetros $(k_0,m_0,0,1)$ pertence a
$\Library_H$, logo o catálogo contém ao menos a árvore-folha e o mínimo é
finito. Todos os conjuntos examinados são finitos. A chave
\eqref{eq:tie-key} é total e depende somente do candidato, não da ordem de
geração, o que prova o determinismo.
\end{proof}

\begin{corollary}[compatibilidade com a enumeração de formas]
\label{cor:shape-enumeration}
Se $\Trees_D(K)$ é o conjunto finito de todas as formas de árvore de
profundidade no máximo $D$, então
\begin{equation}
 C_{D,H}^{\star}(K)
 =\min_{T\in\Trees_D(K)}
   \min_{u,v\in\Values_{D,H}(K)}C_T(K;u,v).
 \label{eq:shape-equivalence}
\end{equation}
O mesmo fluxo canônico vence nos dois lados.
\end{corollary}

\begin{proof}
As classes $\Representations_H(T)$, $T\in\Trees_D(K)$, particionam as
representações pela forma da árvore. Aplique os
\cref{thm:fixed-correctness,thm:variable-correctness} e o mesmo desempate
total.
\end{proof}

\begin{algorithm}[seletor exato por programação dinâmica]
\label{alg:variable-tree}
Construa as tabelas em pós-ordem na árvore diádica completa de profundidade
$D$. Inicialize cada estado com o candidato folha, quando admissível. Una as
tabelas dos filhos pelo valor intermediário $w$, acrescente o marcador e a
penalidade de subdivisão e substitua a entrada somente quando a chave completa
diminuir. A tabela da raiz e os apontadores guardados fornecem o ótimo global
e sua representação.
\end{algorithm}
